\begin{abstract}
% Brief outline of research questions, results.
% (The preferred size of an abstract is one paragraph or one page of text.)
Large language models (LLMs) are increasingly deployed in automated content pipelines, yet remain prone to hallucination, particularly when processing niche, localized information about entities underrepresented in their training data. This thesis addresses hallucination detection in one-sentence LLM-generated summaries of local music venue events, a domain where popularity bias renders both generative models and "LLM-as-a-judge" evaluation approaches unreliable.

We propose and evaluate a hybrid two-phase evaluation architecture that decouples reading comprehension from world knowledge verification. In the first phase, a cross-lingual Natural Language Inference model (XLM-RoBERTa) evaluates semantic entailment between English atomic claims, extracted via a deterministic spaCy dependency parser, and their Dutch source texts, detecting intrinsic hallucinations without an intermediate translation step. In the second phase, an external knowledge graph (MusicBrainz) provides deterministic, non-parametric overrides for extrinsic hallucinations that the NLI model alone cannot resolve.

Evaluated on 47 newsletter summaries from the Nijmegen venue Doornroosje, the baseline NLI pipeline achieved a sentence-level F1-score of 0.8615 and a perfect error detection rate of 1.0000, though at the cost of a high false positive rate of 0.4737. Knowledge graph integration successfully reclassified 5 atomic claims from extrinsic to intrinsic contradictions with 60\% accuracy, demonstrating particular strength in verifying geographic origin claims. However, due to strict min-pooling aggregation, these atomic corrections did not propagate to sentence-level improvements.

These findings demonstrate that a deterministic hybrid NLI and knowledge graph architecture is a methodologically sound alternative to circular LLM-based evaluation for niche, localized entity information, while identifying knowledge graph scope and aggregation strategy as the primary bottlenecks for future improvement.

The full Github Repository can be found at \url{https://github.com/rkoelewijn/bachelor-thesis}
\end{abstract}
